
\documentclass[12pt,a4paper]{article}
\usepackage[utf8]{inputenc}
\usepackage[spanish]{babel}
\usepackage{geometry}
\geometry{margin=2.5cm}
\usepackage{amsmath,amssymb}
\usepackage{graphicx}
\usepackage{hyperref}

\title{\textbf{Informe Académico Automatizado: Proceso KDD y Aprendizaje Estadístico en R}}
\author{Sistema Local RAG - IDEAD}
\date{\today}

\begin{document}

\maketitle

\begin{abstract}
Este reporte analiza la integración del proceso KDD utilizando lenguajes estadísticos como R.
\end{abstract}

\section{Introducción y Contexto}
Este documento ha sido generado automáticamente como parte de la recopilación bibliográfica sobre \textbf{Minería de Datos y KDD}, integrando motores de búsqueda académica y bases de datos vectoriales locales.

\section{Desarrollo Teórico}
El proceso de Knowledge Discovery in Databases (KDD) consta de fases iterativas que van desde la selección de datos hasta la interpretación de patrones...

\section{Referencias}
\begin{enumerate}
    \item Recopilado mediante cadenas booleanas optimizadas para investigación en ingeniería.
\end{enumerate}

\end{document}
